\chapter{Notes: C without a standard library}
\label{ch:freestanding}

Reference notes for the constraints that apply to \lstinline|format/| and
\lstinline|core/|, gathered while designing the on-disk header. Nothing here is
specific to bitterfs; it is the C-language and toolchain background that a first
kernel-level project runs into immediately.

\section{Two conformance modes}
\label{sec:freestanding-hosted}

The C standard defines two kinds of implementation.

\begin{description}[leftmargin=2.5cm,style=nextline]
  \item[Hosted] The normal case. The whole standard library is available, the
    program starts at \lstinline|main|, and the runtime sets up the environment
    before any of your code runs.
  \item[Freestanding] What a kernel, a bootloader, or firmware gets. There is no
    \lstinline|main|, no runtime start-up, and almost no library --- because
    there is no operating system underneath to provide one.
\end{description}

A loadable kernel module is freestanding: it is linked into a running kernel
rather than against libc. That is why \lstinline|printf|, \lstinline|malloc| and
\lstinline|assert| do not exist there, and why the kernel provides its own
\lstinline|printk| and \lstinline|kmalloc| instead.

\section{The freestanding headers}
\label{sec:freestanding-headers}

Freestanding is not \emph{no} headers. The standard requires a small set even in
freestanding mode, and --- the part that is easy to miss --- these ship with the
\textbf{compiler}, not with libc:

\begin{center}
\begin{tabular}{ll}
  \toprule
  Header & Provides \\
  \midrule
  \lstinline|<stdint.h>|  & \lstinline|uint64_t|, \lstinline|UINT64_MAX|, \lstinline|INT64_C()| \\
  \lstinline|<stddef.h>|  & \lstinline|size_t|, \lstinline|ptrdiff_t|, \lstinline|offsetof|, \lstinline|NULL| \\
  \lstinline|<stdarg.h>|  & \lstinline|va_list| and friends \\
  \lstinline|<limits.h>|  & \lstinline|CHAR_BIT|, \lstinline|INT_MAX| \\
  \lstinline|<stdbool.h>| & \lstinline|bool|, \lstinline|true|, \lstinline|false| \\
  \lstinline|<float.h>|, \lstinline|<iso646.h>| & rarely relevant here \\
  \bottomrule
\end{tabular}
\end{center}

They contain no code --- only types, limits and macros --- which is exactly why
they can exist without an operating system.

Their location is worth knowing, because it is what the kernel build removes:

\begin{lstlisting}[style=shell]
$ gcc -print-file-name=include
/usr/lib/gcc/x86_64-linux-gnu/13/include
\end{lstlisting}

GCC's own \lstinline|stdint.h| in that directory is a wrapper that branches on
whether the build is hosted:

\begin{lstlisting}[style=c]
#if __STDC_HOSTED__
# include_next <stdint.h>     /* hand off to glibc's /usr/include/stdint.h */
#else
# include "stdint-gcc.h"      /* the compiler's own self-contained version */
#endif
\end{lstlisting}

\lstinline|-ffreestanding| sets \lstinline|__STDC_HOSTED__| to 0, so it takes
the second branch and needs no libc at all. One file, both worlds.

\section{Why the kernel is the exception}
\label{sec:kernel-nostdinc}

That branch never runs in a kernel module build, because the file is never
found. The kernel's top-level Makefile contains:

\begin{lstlisting}[style=shell]
NOSTDINC_FLAGS += -nostdinc
\end{lstlisting}

and, in 6.12, \textbf{nothing puts the compiler's include directory back}.
Older kernels appended
\lstinline|-isystem $(shell $(CC) -print-file-name=include)| on the same line;
that was removed. A real module build confirms it:

\begin{lstlisting}[style=shell]
hello.c:35:10: fatal error: stdint.h: No such file or directory
\end{lstlisting}

Corroborating evidence: the kernel ships its own
\lstinline|include/linux/stdarg.h|. It would have no reason to carry a copy of
a compiler-provided header unless it had deliberately cut the compiler's
include path.

So the availability of \lstinline|<stdint.h>| is not uniform across this
project:

\begin{center}
\begin{tabular}{lll}
  \toprule
  Directory & Compiled into & \lstinline|<stdint.h>| \\
  \midrule
  \lstinline|format/| & everything          & no \\
  \lstinline|core/|   & the kernel module   & no \\
  \lstinline|user/|   & ordinary programs   & yes \\
  \lstinline|tests/|  & ordinary programs   & yes \\
  \bottomrule
\end{tabular}
\end{center}

Since \lstinline|format/| is included by all four, it takes the strictest
constraint: \textbf{it includes nothing, and defines the integer vocabulary
everything else uses.}

\section{Compiler type macros}
\label{sec:type-macros}

GCC and Clang predefine \lstinline|__UINT64_TYPE__|, \lstinline|__UINT32_TYPE__|
and so on. These look like a lower-level alternative to \lstinline|<stdint.h>|,
but the relationship runs the other way: \textbf{they exist for the header's
benefit}, and came later than it.

\lstinline|<stdint.h>| is mandated by C99, so somebody must write it, and that
somebody needs to know the widths the target actually uses. Historically each
libc shipped its own version full of \lstinline|#if defined(__x86_64__)| chains,
re-deriving what the compiler already knew --- fragile, and capable of
disagreeing with the compiler. The macros let the header simply ask. Clang's
\lstinline|stdint.h| is now largely \lstinline|typedef __UINT64_TYPE__ uint64_t;|.

The header still earns its place, because the typedefs are the small part of it:
it also supplies \lstinline|UINT64_MAX|, the \lstinline|INT64_C()| literal
suffixes, \lstinline|intptr_t|, the \lstinline|int_least*_t| and
\lstinline|int_fast*_t| families, and (via \lstinline|<inttypes.h>|) the
\lstinline|PRIu64| format macros. And \lstinline|uint64_t| is portable in a way
the macros are not --- MSVC defines no such thing.

\section{Three vocabularies for the same integers}
\label{sec:integer-vocabularies}

The same 64-bit unsigned value has three spellings depending on which world the
code lives in, and \lstinline|format/| is included from all of them.

\begin{center}
\begin{tabular}{lll}
  \toprule
  World & Spelling & Source \\
  \midrule
  Standard C          & \lstinline|uint64_t|  & \lstinline|<stdint.h>| \\
  Linux kernel        & \lstinline|u64|, \lstinline|__le64| & \lstinline|<linux/types.h>| \\
  Compiler intrinsic  & \lstinline|__UINT64_TYPE__| & predefined macro \\
  \bottomrule
\end{tabular}
\end{center}

The kernel's \lstinline|__le64| carries a \lstinline|__bitwise| annotation that
means nothing to a normal compiler but lets \textbf{sparse}, the kernel's static
checker, flag any use of a little-endian value where a host-order one was
expected. That is a real safety net, and it is unavailable in userspace.

Defining our own \lstinline|__le64| in \lstinline|format/| would therefore be
actively harmful: in a kernel build it would collide with, or silently shadow,
the annotated one and discard the checking. Any bitterfs types must use their own
prefix.

\section{Marking endianness}
\label{sec:endianness-marking}

Everything on disk is little-endian, and the compiler will not enforce that. A
plain typedef is only a naming convention:

\begin{lstlisting}[style=c]
typedef uint64_t bt_le64;   /* nothing stops a host-order assignment */
\end{lstlisting}

On a little-endian host the mistake even works, which is the worst possible
failure mode --- it surfaces as a corrupt image on somebody else's machine, or
never, which means the bug is in the code forever.

The alternative is a struct wrapper, which makes the compiler refuse any use that
does not go through an accessor:

\begin{lstlisting}[style=c]
typedef struct { uint64_t v; } bt_le64;
\end{lstlisting}

\decide{Plain typedef or struct wrapper? The wrapper buys compile-time
enforcement in every build, including userspace where sparse is unavailable. It
costs clunkier initialisers and noisier field access. Note that the choice is
\emph{not} about header availability --- \lstinline|<stdint.h>| works in all four
consumers --- it is purely about whether the type system enforces the
convention.}

Either way the byte-swapping accessors belong in \lstinline|core/|, not here:
\lstinline|format/| stays declarations only.

\section{Compile-time assertions}
\label{sec:static-assert}

\lstinline|_Static_assert| is a C11 \emph{keyword}, so it needs no header. (The
lower-case \lstinline|static_assert| is a macro from \lstinline|<assert.h>|, which
is \emph{not} a freestanding header --- prefer the keyword.)

\begin{lstlisting}[style=c]
_Static_assert(sizeof(struct bitter_key) == 17, "key size changed");
_Static_assert(__builtin_offsetof(struct bitter_key, offset) == 9, "");
\end{lstlisting}

Assert field offsets as well as total sizes: a size assertion alone will not
catch a reordering that happens to preserve the total. \lstinline|offsetof|
proper comes from \lstinline|<stddef.h>|; \lstinline|__builtin_offsetof| is the
compiler builtin underneath it and needs nothing.

These assertions are the only machine-checked link between this document and the
code. Everything else in the format specification is kept true by hand.

\section{Packing}
\label{sec:packing}

\lstinline|__attribute__((__packed__))| is a GCC and Clang extension, not
standard C. It removes all padding, which is what an on-disk format requires ---
padding must be explicit and documented, never inserted by a compiler whose
choices vary by target.

Two consequences that will surprise you:

\begin{enumerate}
  \item Member accesses become \textbf{unaligned}. Harmless and only slightly
    slower on x86; a fault on some architectures.
  \item \textbf{Taking the address of a packed member warns} under
    \lstinline|-Waddress-of-packed-member|, because the resulting pointer may be
    misaligned and dereferencing it is undefined. Expect this the first time you
    write something like \lstinline|&item->key|; the fix is to copy the member
    out rather than point into the structure.
\end{enumerate}

\section{Verifying layout with pahole}
\label{sec:pahole}

\lstinline|pahole| (from the \lstinline|dwarves| package, already installed for
the kernel build) reads DWARF debug information and prints a structure's real
layout --- including \textbf{holes} where padding was inserted.

\begin{lstlisting}[style=shell]
echo '#include "format/bitterfs_format.h"' > /tmp/t.c
gcc -g -c -std=c11 -Wall -Wextra -I. /tmp/t.c -o /tmp/t.o
pahole /tmp/t.o
\end{lstlisting}

A hole you did not declare means the structure is wrong. The output is also,
near enough, the field table this document needs --- read the offsets off it
rather than computing them by hand, and the two cannot disagree.

\section{Accessors the optimiser can see through}
\label{sec:endian-accessors}

Everything on disk is little-endian, and \lstinline|core/bitter_endian.h| is
the only file that knows it. The accessors there are the most frequently
executed code in the filesystem --- a B-tree search binary-searches roughly 160
items per node across several levels, and every key comparison reads three
fields --- so their shape was settled by measurement rather than by taste. The
measurements are recorded here because the resulting code looks worse than the
obvious version and there would otherwise be a standing temptation to tidy it.

All figures are gcc 13, \lstinline|-O2|, x86-64.

\subsection{The signature}

\begin{lstlisting}[style=c]
static inline bt_u64 bt_get_le64(const void *p);
static inline void   bt_put_le64(void *p, bt_u64 v);
\end{lstlisting}

The little-endian side is always behind a \textbf{pointer} and the host side is
always a plain value. That is not a stylistic choice: byte order is a property
of a layout in memory, and a value sitting in a register does not have one.
A signature like \lstinline|bt_get_le64(bt_le64 v)| would be meaningless.

Two further points about \lstinline|const void *| rather than
\lstinline|const bt_le64 *|:

\begin{enumerate}
  \item Every call site passes the address of a \emph{packed} member.
    Measured: passing \lstinline|&it->key.offset| to a
    \lstinline|const bt_u64 *| parameter warns under
    \lstinline|-Waddress-of-packed-member|; passing it to a
    \lstinline|const void *| parameter does not.
  \item \lstinline|bt_le64| and \lstinline|bt_u64| are both
    \lstinline|typedef __UINT64_TYPE__| --- the \emph{same type} to the
    compiler (\cref{sec:endianness-marking}). A typed parameter would appear
    to check something while checking nothing. \lstinline|const void *| states
    what is true: raw bytes, interpreted here.
\end{enumerate}

The getter's local must be \lstinline|const unsigned char *|. Dropping the
qualifier compiles, but emits
\lstinline|initialization discards 'const' qualifier| under
\lstinline|-Wdiscarded-qualifiers|, and silencing that with a cast rather than
a qualifier discards the only enforcement the signature had. It matters in
practice: \lstinline|bitter-dump| has every reason to \lstinline|mmap| an image
\lstinline|PROT_READ|, and writing through a stripped-\lstinline|const| pointer
to a genuinely read-only object is a fault, not a wrong answer.

\subsection{Bytes, not a pointer cast}

The tempting implementation is \lstinline|return *(const bt_u64 *)p;|. It is
wrong twice over:

\begin{description}[leftmargin=2.2cm,style=nextline]
  \item[Alignment] The structures are packed (\cref{sec:packing}), so many of
    these addresses are unaligned. Dereferencing a misaligned
    \lstinline|bt_u64 *| is undefined --- tolerated on x86, a fault elsewhere.
  \item[Aliasing] Character types are the sole exception to strict aliasing:
    reading \emph{any} object through an \lstinline|unsigned char *| is
    explicitly permitted. That rule is what makes the byte-wise version defined
    rather than merely working today.
\end{description}

Note which thing is cast. The bytes are indexed first and the resulting
\textbf{value} is widened; the \textbf{pointer} is never cast. C's precedence
gives this for free --- \lstinline|[]| is postfix and binds tighter than a
cast, so \lstinline|(bt_u64)b[1]| parses as \lstinline|(bt_u64)(b[1])|.
Casting the pointer requires going out of one's way to write
\lstinline|((const bt_u64 *)b)[1]|.

\subsection{The promotion trap}

The cast must sit inside the shift, not around it:

\begin{center}
\begin{tabular}{ll}
  \toprule
  Expression & Result for \lstinline|b[7] == 0x80| \\
  \midrule
  \lstinline|(bt_u64)b[7] << 56|   & \lstinline|0x8000000000000000| \\
  \lstinline|(bt_u64)(b[7] << 56)| & \lstinline|0x0000000000000000| \\
  \bottomrule
\end{tabular}
\end{center}

\lstinline|b[7]| promotes to \lstinline|int|, and shifting a 32-bit
\lstinline|int| by 56 is undefined behaviour --- the whole top byte disappears.
GCC caught this only because the shift count was a literal constant
(\lstinline|-Wshift-count-overflow|); written as a loop with a variable count,
the warning vanishes and the bug does not. UBSan still reports it at runtime,
which is one concrete reason \lstinline|make -C user SANITIZE=1| exists.

The putter is the mirror image, and the cast moves to the other end:

\begin{lstlisting}[style=c]
b[7] = (unsigned char)(v >> 56);   /* v is already 64-bit; narrow on the way out */
\end{lstlisting}

Getter casts first, putter casts last. If the two end up written the same way,
one of them is wrong.

\subsection{Why the unrolled version, measured}

The natural implementation is a loop. It is correct, and it is what anybody
would write:

\begin{lstlisting}[style=c]
bt_u64 num = 0;
for (int i = 0; i < 8; i++)
    num += (bt_u64)b[i] << 8*i;
return num;
\end{lstlisting}

It is also forty instructions where the alternative is one:

\begin{center}
\begin{tabular}{ll}
  \toprule
  Form & \lstinline|-O2| output \\
  \midrule
  loop, \lstinline|+=|                          & 8-iteration loop, ${\sim}40$ instructions \\
  loop, \lstinline@|=@                          & identical --- the operator is not the cause \\
  loop, \lstinline|+=|, \lstinline|-funroll-loops| & unrolled, still 8 loads $+$ 8 shifts \\
  unrolled expression, \lstinline@|@            & \lstinline|movq (%rdi), %rax| \\
  unrolled expression, \lstinline|+|            & \lstinline|movq (%rdi), %rax| \\
  \bottomrule
\end{tabular}
\end{center}

The last two rows are the informative ones: the operator is irrelevant, so the
\textbf{loop} is what blocks the optimisation. GCC's byte-combining idiom
recogniser matches a straight-line expression tree and does not reconstruct one
from a loop --- not even after forced unrolling, because by that point the
shift counts originate in an induction variable rather than being constants in
the source.

So the accessors are written out longhand, all eight terms, and the loop is
preserved as a comment in \lstinline|bitter_endian.h| so the reasoning is not
rediscovered later. Use \lstinline@|@ rather than \lstinline|+| in the surviving
version: the compiler emits the same instruction either way, but these are
disjoint bit ranges being assembled rather than numbers being summed, and a
reader who sees \lstinline|+| must stop and confirm that no term can carry into
the next.

\subsection{What this buys on a big-endian host}

The same source compiles to a load plus a \lstinline|bswap|. That is the whole
reason the file contains no \lstinline|#ifdef| for byte order --- and, in turn,
why \lstinline|core/| stays free of environment branches and can be compiled
unchanged into both the kernel module and the userspace tools.
